\textbf{Phần 1. Trắc nghiệm lựa chọn}
\setcounterpageref{ex}{0}

\begin{ex} %[3D2V77]
Tính giá trị phép tính: $45\,280 + 38\,450$.
\choice
{$83\,630$}
{$83\,830$}
{\True $83\,730$}
{$82\,730$}
\loigiai{
$45\,280 + 38\,450 = 83\,730$.
\textbf{Đáp án của câu hỏi là C}
}
\end{ex}

\begin{ex} %[3D2V77]
Tính giá trị phép tính: $92\,400 - 45\,800$.
\choice
{\True $46\,600$}
{$47\,600$}
{$46\,400$}
{$56\,600$}
\loigiai{
$92\,400 - 45\,800 = 46\,600$.
\textbf{Đáp án của câu hỏi là A}
}
\end{ex}

\begin{ex} %[3D2V77]
Tìm số tự nhiên $x$, biết: $x + 28\,500 = 75\,000$.
\choice
{$47\,500$}
{$46\,000$}
{$45\,500$}
{\True $46\,500$}
\loigiai{
$x = 75\,000 - 28\,500 = 46\,500$.
\textbf{Đáp án của câu hỏi là D}
}
\end{ex}

\begin{ex} %[3D2V77]
Tính nhanh giá trị: $34\,500 + 28\,000 + 15\,500 + 12\,000$.
\choice
{$88\,000$}
{\True $90\,000$}
{$92\,000$}
{$89\,000$}
\loigiai{
$(34\,500 + 15\,500) + (28\,000 + 12\,000) = 50\,000 + 40\,000 = 90\,000$.
\textbf{Đáp án của câu hỏi là B}
}
\end{ex}

\begin{ex} %[3D2V77]
Một kho gạo có $85\,000\text{ kg}$ gạo. Người ta xuất đi hai đợt: đợt một $28\,000\text{ kg}$, đợt hai $32\,500\text{ kg}$. Trong kho còn lại bao nhiêu ki-lô-gam gạo?
\choice
{$25\,500\text{ kg}$}
{$24\,000\text{ kg}$}
{\True $24\,500\text{ kg}$}
{$34\,500\text{ kg}$}
\loigiai{
Tổng xuất: $28\,000 + 32\,500 = 60\,500\text{ kg}$. Còn lại: $85\,000 - 60\,500 = 24\,500\text{ (kg)}$.
\textbf{Đáp án của câu hỏi là C}
}
\end{ex}

\begin{ex} %[3D2V77]
Một công ty quý 1 thu được $45\,000\,000$ đồng, quý 2 thu được nhiều hơn quý 1 là $12\,000\,000$ đồng. Cả hai quý công ty thu được bao nhiêu tiền?
\choice
{\True $102\,000\,000$ đồng}
{$90\,000\,000$ đồng}
{$100\,000\,000$ đồng}
{$95\,000\,000$ đồng}
\loigiai{
Quý 2: $45 + 12 = 57$ triệu. Cả hai quý: $45 + 57 = 102$ triệu đồng.
\textbf{Đáp án của câu hỏi là A}
}
\end{ex}

\begin{ex} %[3D2V77]
Tìm hiệu của số lớn nhất có 5 chữ số và số bé nhất có 5 chữ số.
\choice
{$90\,000$}
{$99\,999$}
{$88\,888$}
{\True $89\,999$}
\loigiai{
$99\,999 - 10\,000 = 89\,999$.
\textbf{Đáp án của câu hỏi là D}
}
\end{ex}

\begin{ex} %[3D2V77]
Tính giá trị biểu thức: $100\,000 - (35\,000 + 45\,000)$.
\choice
{$30\,000$}
{\True $20\,000$}
{$25\,000$}
{$15\,000$}
\loigiai{
$100\,000 - 80\,000 = 20\,000$.
\textbf{Đáp án của câu hỏi là B}
}
\end{ex}

\begin{ex} %[3D2V77]
Tổng hai số là 84 000, hiệu hai số là 16 000. Số lớn là
\choice
{$42\,000$}
{$34\,000$}
{\True $50\,000$}
{$48\,000$}
\loigiai{
Số lớn: $(84\,000 + 16\,000) : 2 = 50\,000$.
\textbf{Đáp án của câu hỏi là C}
}
\end{ex}

\begin{ex} %[3D2V77]
Một mảnh vườn hình chữ nhật có chu vi $600\text{ m}$, chiều dài hơn chiều rộng $60\text{ m}$. Chiều dài của mảnh vườn là
\choice
{\True $180\text{ m}$}
{$120\text{ m}$}
{$150\text{ m}$}
{$200\text{ m}$}
\loigiai{
Nửa chu vi: $600 : 2 = 300\text{ m}$. Chiều dài: $(300 + 60) : 2 = 180\text{ (m)}$.
\textbf{Đáp án của câu hỏi là A}
}
\end{ex}

\begin{ex} %[3D2V77]
Nếu số bị trừ tăng 5000 đơn vị và số trừ giữ nguyên thì hiệu sẽ
\choice
{Giảm đi 5000 đơn vị}
{Không đổi}
{Tăng thêm 10 000 đơn vị}
{\True Tăng thêm 5000 đơn vị}
\loigiai{
Số bị trừ tăng bao nhiêu thì hiệu tăng bấy nhiêu.
\textbf{Đáp án của câu hỏi là D}
}
\end{ex}

\begin{ex} %[3D2V77]
Một cuộn dây cáp dài $50\,000\text{ m}$, đã dùng hết $32\,400\text{ m}$. Cuộn dây cáp còn lại dài bao nhiêu mét?
\choice
{$18\,600\text{ m}$}
{\True $17\,600\text{ m}$}
{$17\,400\text{ m}$}
{$16\,600\text{ m}$}
\loigiai{
$50\,000 - 32\,400 = 17\,600\text{ (m)}$.
\textbf{Đáp án của câu hỏi là B}
}
\end{ex}

\textbf{Phần 2. Trắc nghiệm đúng sai}
\setcounterpageref{ex}{0}

\begin{ex} %[3D2V77]
Xét tính đúng sai của các phép toán cộng trừ phạm vi 100 000:
\choiceTF
{\True Phép cộng và trừ được thực hiện lần lượt từ phải sang trái.}
{$60\,000 - 25\,000 = 45\,000$.}
{\True Khi cùng thêm vào cả hai số hạng một số đơn vị thì tổng tăng thêm gấp đôi số đơn vị đó.}
{$75\,000 + 25\,000 = 90\,000$.}
\loigiai{
\begin{itemize}
    \item Ý a) Đúng: Quy tắc cơ bản.
    \item Ý b) Sai: Sai vì $60\,000 - 25\,000 = 35\,000$.
    \item Ý c) Đúng: Tính chất tổng.
    \item Ý d) Sai: Sai vì $75\,000 + 25\,000 = 100\,000$.
\end{itemize}
\begin{itemize}
    \item a) \textbf{Đúng}.
    \item b) \textbf{Sai}.
    \item c) \textbf{Đúng}.
    \item d) \textbf{Sai}.
\end{itemize}
}
\end{ex}

\begin{ex} %[3D2V77]
Một kho dự trữ có 90 000 bao xi măng:
\choiceTF
{\True Xuất đi 45 000 bao thì còn lại đúng một nửa.}
{\True Nếu nhập thêm 10 000 bao nữa thì kho có đúng 100 000 bao xi măng.}
{Xuất đi 60 000 bao thì số bao còn lại nhiều hơn số bao đã xuất.}
{\True Tổng số bao nhập và xuất trong ngày có thể tính bằng phép cộng.}
\loigiai{
\begin{itemize}
    \item Ý a) Đúng: $45\,000 : 90\,000 = \dfrac{1}{2}$.
    \item Ý b) Đúng: $90\,000 + 10\,000 = 100\,000$ bao.
    \item Ý c) Sai: Sai vì còn 30 000 bao, ít hơn 60 000 bao.
    \item Ý d) Đúng: Đúng.
\end{itemize}
\begin{itemize}
    \item a) \textbf{Đúng}.
    \item b) \textbf{Đúng}.
    \item c) \textbf{Sai}.
    \item d) \textbf{Đúng}.
\end{itemize}
}
\end{ex}

\textbf{Phần 3. Trả lời ngắn}
\setcounterpageref{ex}{0}

\begin{ex} %[3D2V77]
Tính: $64\,500 + 28\,700$.
\shortans{93200}
\loigiai{
\textbf{Đáp số:} $93200$\\
$64\,500 + 28\,700 = 93\,200$.
}
\end{ex}

\begin{ex} %[3D2V77]
Tính: $82\,000 - 45\,600$.
\shortans{36400}
\loigiai{
\textbf{Đáp số:} $36400$\\
$82\,000 - 45\,600 = 36\,400$.
}
\end{ex}

\begin{ex} %[3D2V77]
Tìm $x$, biết: $x - 24\,000 = 56\,000$.
\shortans{80000}
\loigiai{
\textbf{Đáp số:} $80000$\\
$x = 56\,000 + 24\,000 = 80\,000$.
}
\end{ex}

\begin{ex} %[3D2V77]
Tính nhanh: $45\,000 + 18\,500 + 15\,000 + 11\,500$.
\shortans{90000}
\loigiai{
\textbf{Đáp số:} $90000$\\
$60\,000 + 30\,000 = 90\,000$.
}
\end{ex}

\begin{ex} %[3D2V77]
Tìm hiệu giữa 100 000 và 72 500.
\shortans{27500}
\loigiai{
\textbf{Đáp số:} $27500$\\
$100\,000 - 72\,500 = 27\,500$.
}
\end{ex}

\begin{ex} %[3D2V77]
Tổng hai số là 70 000, số bé là 30 000. Tìm số lớn.
\shortans{40000}
\loigiai{
\textbf{Đáp số:} $40000$\\
Số lớn là: $70\,000 - 30\,000 = 40\,000$.
}
\end{ex}
